\documentclass[script,pro]{debate}
\motion{「勇敢的人先享受世界」是一个陷阱 / 不是一个陷阱}
\stance{正方}{是一个陷阱}
\meta{赛制}{中国科学技术大学新生辩论赛（正方先立论，正四最后结辩）}
\meta{口径来源}{备赛文档 第十节 10.1 正方口径表。定义、判准、论点、数据均取自该表；允许换说法，不允许换意思、换数字。}
\begin{document}
\debatetitle
\begin{thesisbox}[全场一句话立论]它把「有没有兜底」这张门票，标成了「够不够勇敢」这枚勋章，于是最付不起代价的人被叫去付账。\end{thesisbox}
\begin{thesisbox}[全场只守两句话]① 陷阱不需要有人挖。② 省略条件只是不严谨，拿勇敢去顶替条件才是陷阱。\end{thesisbox}

\begin{speech}{1. 正一开篇立论稿}{正文 839 字 / 180 秒}
谢谢主席。开场二十秒，我先把最要紧的那个词说清楚。

陷阱，词典上有两个意思，一个是捕兽的深坑，一个是使人受骗上当的圈套。两个意思有同一个内核：看着是路，走进去是坑，而且正是它的外观把你引了进去。请评委记住，陷阱不需要有人挖。中等收入陷阱、消费陷阱，都没有设局者。另外两个词按常识讲：勇敢的人，是后果不确定还选择行动的人；先享受世界，是比不敢动的人更早拿到好处。

所以判断一句话是不是陷阱，要看三件事。第一，它说的因果，和真实的因果对不对得上。第二，照着做的人，受损是常态还是例外。第三，这份损失，是不是被它遮住的那部分东西造成的。三件都成立，它就是陷阱。缺一件，它就只是一句不够严谨的话。三件我方全扛，对方打掉一件即可。

我方第一点，门票换勋章。能不能说走就走，本来是一道算得出来的题：存款够几个月，家里要不要你补贴，简历断一年补不补得回来。这句话把这道题换掉了，换成了另一道题，你够不够勇敢。新的这道题没有答案，只能靠感觉回答。你想想，一道没有单位的题，人怎么答得出来？社会心理学家 Ross 一九七七年讲过一个词，叫基本归因错误，说的就是人总高估品质，低估处境。智联招聘二〇二五年春季那份三万七千多人的调研里，两成四的人明说不会 gap 三个月，理由是经济压力。另有将近两成，怕的是简历断档。四成多人算过账，然后留下了。媒体早就记下了它的变化：勇敢，从说走就走的冲动，变成了有说走就走的物质实力。门槛一直都在，它只是改了个名字。请对方回答，一个把「我算过账，我走不了」翻译成「你不够勇敢」的说法，对这四成人是什么？

我方第二点，删过的账本。我们听到的所有「我勇敢了，我过得很好」，都来自飞回来的那一架飞机。一九四三年 Wald 研究战机装甲时就发现，只看活着回来的样本，结论会被推反。这句话最主要的说服力还有一句，你以后会后悔没去。可是二〇二三年，一份近两千六百人的公开复制做出来的结果是，回顾一生，五成一的人更困扰于做过的事，四成九更困扰于没做的事。也就是说，连这句话最常用的那个理由，在最好的证据下都只是五五开。它卖的不是亏本生意，\stress{是一张写着稳赚的账单。}

我方全场只反复问一个问题：这句话里，哪个字告诉过你门票多少钱？

最后我方先立一条规矩，免得后面吵不清：省略条件只是不严谨，拿另一个变量去顶替被省略的那个，才是陷阱。谢谢。
\end{speech}
\begin{contingency}
\ifthen{对方说「陷阱必须有设局者」}{立刻回：「那中等收入陷阱是谁设的局？请对方不要用定义取消今天的比赛。」}
\ifthen{对方问「那什么不是陷阱」}{三秒内答：「说了条件的经验分享不是，只省略、没顶替的格言不是。」}
\ifthen{时间不够}{删 Wald 那一句的年份与背景，保留「飞回来的那一架飞机」。}
\end{contingency}

\begin{stage}{2. 正一被质询防守预案}{反四质询，120 秒双边计时}
\textbf{原则}：双边计时，每答控制在八秒内。关键前提不含糊，用「是，但前提是」。不否认常识。被逼二选一时，先指出预设有问题，再给我方的第三种说法。
\begin{qa}
\qarow{陷阱是不是必须有人设局？}{不是。中等收入陷阱没有设局者，它照样叫陷阱。}
\qarow{那按你方标准，什么不是陷阱？}{说了条件的经验分享不是；只省略、没顶替的格言不是。}
\qarow{「早睡早起身体好」是不是陷阱？}{不是。它没有告诉你，睡不好是因为你不自律。}
\qarow{这句话有没有承诺你一定过得好？}{没有承诺结果。它承诺的是顺序，而顺序也是承诺。}
\qarow{你方能不能拿出照它做的人普遍变差的数据？}{没有直接数据，这句话才流行三年。我方打的是代价错配。}
\qarow{那你方第二件事不就没立住？}{立住了。受损的常态就是付不起的人按付得起的人的活法行动。}
\qarow{你方那四成多人不是没走吗？他们受什么损了？}{他们被这句话判成了不勇敢。这正是外观与真实因果不一致。}
\qarow{勇敢和鲁莽你方分不分？}{分。可这九个字里没有一个字提到后果，分不出来的是它。}
\qarow{是不是这句话让人不算账的？}{我方不主张每个人都不算账。我方主张它把能算的题换成了不能算的题。}
\qarow{你方是不是在劝人别勇敢？}{不是。我方劝的是先看清门票价格，这和勇敢不冲突。}
\end{qa}
\end{stage}

\begin{stage}{3. 正四质询反一问题链}{120 秒，双边计时}
\textbf{总目标}：拿到「陷阱不需要设局者」这个承认，并逼反方正面说出「这句话到底承诺了什么」。
\begin{chain}{链一：设局者链}{堵死反方的定义杀}
\q{中等收入陷阱有设局者吗？}{答「没有」：进 2；答「有」：问「是谁」}
\q{那它算不算陷阱？}{答「算」：进 3；答「不算」：记下来，小结重锤}
\q{所以陷阱可以没有人挖，对吗？}{任何回答都进链二}
\end{chain}
\begin{chain}{链二：承诺链}{逼出正面回答}
\q{这句话承诺过你会过得更好吗？}{答「没有」：进 2；答「有」：直接收}
\q{那它凭什么能让人辞职？}{答「它不能」：进 3；答「靠共鸣」：问「共鸣什么」}
\q{一句什么都没说的话，能红三年吗？}{不等答完，进链三}
\end{chain}
\begin{chain}{链三：门票链}{拿到「条件决定能不能走」的承认}
\q{存款够不够，影响不影响能不能裸辞？}{答「影响」：进 2；答「不影响」：报智联数据}
\q{这九个字里，哪个字提到了存款？}{任何回答都进 3}
\q{所以决定的是条件，说出口的是勇敢，对吗？}{收}
\end{chain}
\cut{好的，我理解为您没有否认。下一个问题。}
\close{感谢。对方已经承认陷阱不需要设局者，也没能说出这句话承诺了什么，我方二辩会继续。}
\end{stage}

\begin{speech}{4. 正二驳论稿}{正文 438 字 + 临场位 99 字 / 120 秒}
对方立论最大的问题只有一个：\stress{他们把这九个字改写了一遍，然后来守那个改写过的版本。}

第一，对方说它给的是顺序不是支票，说它承诺的只是「先行动的人更早开始」。可请评委再听一遍原句，是「勇敢的人先享受世界」，不是「勇敢的人先开始」。享受两个字，是对方自己删掉的。一句话如果真的只承诺了「开始的人开始了」，那它就是一句同义反复，请问同义反复怎么可能让人辞职，又怎么可能被骂三年？而且对方自己也说了，它推荐先动。推荐里带着一个先字，先就是排序，排序就是在告诉你，不这么做的人排在后面。

第二，对方说摔跤的是鲁莽不是勇敢。这个注脚听上去很合理，可它是对方加的。我方请对方指出来，这九个字里，哪一个字提醒你去算后果？指不出来，对方今天守的就不是这道题，而是另一句话：准备好的人先享受世界。那正是我方在说的东西。

第三，对方引一九九四年的后悔研究，说人更后悔没做的事。我方承认这项研究存在，但请对方把二〇二三年那份近两千六百人的公开复制一起念完：长期回顾是五成一对四成九，那条结论没有复制出来。对方引的，是三十年前一个小样本上被推翻的那一半。

\reserve{约 100 字，回应反二刚才最狠的一句。若反二打「你方拿不出照它做的人普遍变差的数据」，回：没有直接数据，因为这句话才流行三年；可对方同样拿不出「他们都过得很好」。双方都在打机制，那就请评委看机制：谁的机制里，付不起的人被叫去付账。}

回到判准：外观与真实因果对不上，这一件对方今天一个字都没碰。谢谢。
\end{speech}
\begin{contingency}
\ifthen{反二没有打我方第二点}{把第三驳点压缩到 30 字，把时间给第一驳点，追问「享受两个字去哪了」。}
\ifthen{反二承认了陷阱不需要设局者}{开头改成：「感谢对方二辩，设局者这一格我们已经没有分歧了。」}
\end{contingency}

\begin{stage}{5. 正二对辩要点}{90 秒，单边计时}
\begin{points}{进攻点（4–5 个，每个 15–20 秒，说完就停）}
\item 享受两个字是您方删的，请问它去哪了？
\item 这九个字里哪个字提醒您去算后果？请给一个字。
\item 您方承认陷阱不需要设局者了吗？一辩说不需要，您同意吗？
\item 后悔研究请念完二〇二三年那一半，五成一对四成九。
\item 四成多人算完账留下了，您方管这叫「他们没被骗」，我方管这叫「他们被判成了不勇敢」。
\end{points}
\begin{points}{备用点（6–8 个，用于回应）}
\item 省略是不严谨，顶替才是陷阱，这条线我方全场不改。
\item 校友分享会会说「我当时家里能兜」，这句话不说。
\item 中等收入陷阱没有设局者，您方不能又用它又否认它。
\item 您方的三成五裸辞，测的是意愿，不是结果。
\item 不含在校生的青年失业率还在一成五左右，空窗不是中性的。
\item 我方从不说勇敢不好，我方说的是门票被改了名。
\item 您方越强调要准备，越是在承认门槛存在。
\item 一句不影响任何人的话，不会被骂三年。
\end{points}
\begin{points}{对方最可能问的三个问题 → 我方回应}
\item 「照它做的人普遍变差，数据呢？」→ 没有直接数据，这句话才流行三年；您方也拿不出反面数据。双方都在机制层。
\item 「那什么不是陷阱？」→ 说了条件的经验分享不是，只省略没顶替的格言不是。
\item 「你方是不是在劝人别勇敢？」→ 不是。先看清门票价格，然后再决定买不买。
\end{points}
\end{stage}

\begin{stage}{6. 正三盘问问题链}{90 秒，单边计时}
\textbf{赛前确认}：向组委会确认被盘问方回答是否计时；若不计时，每条链可多加一问。
\begin{chain}{链一：口径链（先问反二，再问反四）}{检验两人对设局者的答案是否一致}
\q{（问反二）陷阱需不需要有设局者？}{答「不需要」：进 2；答「需要」：记下来}
\q{（问反四）四辩，同样的问题，您答什么？}{两人不一致就当场点名，小结重锤}
\q{（问反二）那您方一辩说的「没有收款方」，还算不算理由？}{答「算」：指出与刚才矛盾}
\end{chain}
\begin{chain}{链二：承诺链}{逼反方在「享受」两个字上表态}
\q{（问反四）这句话里的「享受」，承诺了什么？}{答「什么都没承诺」：进 2}
\q{（问反四）那把享受换成开始，意思变不变？}{答「不变」：进 3；答「变」：收}
\q{（问反二）二辩，您也认为不变吗？}{两人答案不同即收束}
\end{chain}
\textbf{盘问目标}：小结里要证明两件事——对方在设局者这一格已经让步；对方为了守住立论，必须把原句改写成一句同义反复。
\end{stage}

\begin{stage}{7. 正二、正四被盘问防守预案}{两人口径必须一致}
\begin{qaa}
\qaarow{你方能不能证明受损是常态？}{能。常态是代价错配：付不起的人按付得起的人的活法行动。}{能。常态是代价错配，这一点我方二辩已经说过。}
\qaarow{有没有直接数据？}{没有直接数据，这句话才流行三年。}{没有。对方同样拿不出「他们都过得很好」。}
\qaarow{那什么不是陷阱？}{说了条件的经验分享不是。}{只省略、没顶替的格言不是。}
\qaarow{你方承不承认这句话有积极作用？}{承认有。有积极作用的东西也可以是陷阱。}{承认。我方打的不是好坏，是性质。}
\qaarow{勇敢和鲁莽你方分不分？}{分。分不出来的是这九个字。}{分。请对方指出哪个字提醒了算后果。}
\qaarow{你方是不是在替躺平说话？}{不是。我方说的是先看门票价格。}{不是。看清价格再买，和勇敢不冲突。}
\end{qaa}
\end{stage}

\begin{speech}{8. 正三盘问小结稿}{正文 398 字 / 90 秒}
刚才的盘问里，对方交出了两样东西。

第一样，对方承认陷阱不需要有设局者。那么对方一辩稿里那句「没有签名，也没有收款方」，就再也不能当作理由了。今天这道题不是在问谁挖的坑，是在问那层浮土在不在。

第二样，对方为了守住自己的立论，把「享受」换成了「开始」。各位评委，这是一次改写。原句是勇敢的人先享受世界。对方的版本是勇敢的人先开始。开始的人开始了，这是一句什么都没说的话。可一句什么都没说的话，不会让人辞职，也不会被骂三年。对方只有两条路：承认它说了「享受」，那就要交出享受的门票价格；或者坚持它没说，那就请对方解释它凭什么流行。

说到底，双方真正的分歧只剩一格：这句话有没有把门槛换成勇气。对方全场的回应是「人照样在算账」。可我方从来没有说人不算账。我方说的是，算完账走不了的那四成多人，被这句话判成了不勇敢。这就是外观与真实因果的不一致。一个人算清楚了自己走不了，还要被这句话补上一句「你不够勇敢」，这不是鼓励，这是判决。请评委记住，我方全场问的那个问题，对方到现在没有回答：这句话里，哪个字告诉过你门票多少钱？
\end{speech}

\begin{stage}{9. 自由辩战场设计}{240 秒}
\begin{battleground}{战场一：门票（前 90 秒，主战场）}
\begin{fields}
\field{目标}{让评委看到，决定能不能走的是条件，说出口的却是勇敢。}
\field{开场问题}{这九个字里，哪个字提到了存款？}
\field{对方答「没有，但也没说不要算」→ 追问}{没说不要算，和把问题换成「你敢不敢」，是两回事。请正面回答，它把问题换了没有。}
\field{对方答「条件当然重要，这是常识」→ 追问}{那这句常识，为什么在这九个字里一个字都没有？}
\field{对方可能反攻 → 我方防守}{对方会说「四成多人照样算账」。回：他们算了，然后被判成不勇敢。那四成人不是反例，是受害者。}
\field{收束语}{对方承认条件决定一切，却说不出这句话里哪个字提到条件。门票在，标签换了，这就是我方说的陷阱。}
\end{fields}
\end{battleground}
\begin{battleground}{战场二：账本（中 90 秒）}
\begin{fields}
\field{目标}{让评委看到这句话的说服力建立在一个已经被推翻的常识上。}
\field{开场问题}{后悔研究的二〇二三年那一半，对方念不念？}
\field{对方答「五五开就不是陷阱」→ 追问}{五五开的是后悔，不是代价。请问代价那一边，谁替没兜底的人付？}
\field{对方答「幸存者偏差人人都有」→ 追问}{校友分享会会说家里能兜，这句话说吗？说了条件的叫经验，删了的才叫陷阱。}
\field{对方可能反攻 → 我方防守}{对方会说我方拿不出受损数据。回：承认拿不出直接数据，对方也拿不出反面数据，那就比机制。}
\field{收束语}{对方引的那半句已经被复制推翻，剩下的半句正好是我方的：做和不做，长期看一样。那凭什么只劝人做？}
\end{fields}
\end{battleground}
\begin{points}{必问（前两分钟内问完）}
\item 这九个字里，哪个字提到了存款？
\item 「享受」两个字，对方删到哪里去了？
\item 陷阱需不需要有设局者？请给一个字的回答。
\item 二〇二三年那份复制，对方念不念完？
\end{points}
\begin{points}{必答（15 秒以内正面回答）}
\item 「照它做的人普遍变差，数据呢？」→ 没有直接数据，这句话才流行三年；我方打的是代价错配，对方同样没有反面数据。
\item 「那什么不是陷阱？」→ 说了条件的经验分享不是，只省略、没顶替的格言不是。
\item 「你方是不是在劝人别勇敢？」→ 不是。先看清门票价格，再决定买不买。
\end{points}
\textbf{时间分配与站位}：前 90 秒战场一，抛完四个必问；中 90 秒战场二；最后 60 秒只重复「哪个字提到了存款」，不开新战场，留 20 秒备用。一辩守定义与判准，二辩接驳论过的点，三辩接盘问成果，四辩收束与价值。
\end{stage}

\begin{speech}{10. 正四结辩稿}{正文 771 字 / 210 秒，留 30–40 秒临场回应}
谢谢主席。全场打到现在，分歧其实只剩一格。

对方说，这句话给的是顺序，不是支票，所以没有落差。我方说，落差就在「享受」那两个字里。对方为了把落差抹平，全场做了一件事：把「先享受世界」读成「先开始」。可原句不是这么写的。请评委记住，一个需要被改写才守得住的立论，本身就说明原句有问题。

我要正面处理对方全场最有力的那一点。对方说，四成多人明明算过账，所以没有伪装。我方承认，他们算了。可各位，他们算完之后听到的是什么？是这句话在替他们下结论：走不了的人，是不够勇敢的人。这就是我方说的那次顶替。省略条件只是不严谨，拿勇敢去顶替条件，才是陷阱。这条线我方从一辩立到现在，对方一次也没有正面回应。

我方两个论点，今天都还站着。门票换勋章，对方没有指出这九个字里哪个字提到了条件。删过的账本，对方引的那半句已经被二〇二三年的复制推翻，而剩下的那半句是五成一对四成九，做和不做长期看一样。那我想问，既然一样，这句话凭什么只劝人做，还要加一个先字？

还有一件事今天双方没有分歧：条件决定了谁能走。对方三辩、四辩都承认了。可承认之后，对方没有再往前走一步，去看看这句话是怎么处理这个条件的。它没有隐瞒条件，它做了一件更巧的事，把条件改了个名字，叫勇敢。名字一换，账就算不出来了。勇敢没有单位。

我方全场问了四次的那个问题，对方一次也没有回答：这句话里，哪个字告诉过你门票多少钱？

最后我想说件小事。

晚上十一点。宿舍熄了灯，还有人在看手机。屏幕里是别人辞职之后的极光，配文写着那九个字。他把手机扣过来，心里冒出来的第一个念头不是「我没钱」，是「我不够勇敢」。

这就是今天我方最想让各位看见的那一步。那句话没有骗他的钱，它换掉了他心里那道题。原来的题他答得出来：存款三个月，房租两千，家里每个月还要他寄一点。新的题他答不出来，因为勇敢没有单位，不能算，只能比。

我方从来不反对有人出发。我方反对的是，让一个人把「我付不起」记成「我不敢」。这两句话听起来像，但一句是账单，一句是判决。

今天这道题不是在争该不该勇敢。是在争，谁有资格替一个人定义他的胆量。

我方认为，「勇敢的人先享受世界」是一个陷阱。谢谢。
\end{speech}
\begin{contingency}
\ifthen{反四把我方打成「什么都叫陷阱」}{临场位第一句就用：「我方从一辩就立了线，省略不算，顶替才算。请评委翻回去看。」}
\ifthen{反四主打「叫它陷阱会让人不敢动」}{回：「那是后果，不是性质。而且这句话本身就承认了那里有个坑。」}
\ifthen{反四没有处理『享受』两个字}{临场位直接点名：「反方四辩到最后也没有告诉我们，享受两个字去哪了。」}
\end{contingency}
\end{document}
